\documentclass[10pt]{beamer}

\usepackage{verbatim}
%\usepackage[latin1]{inputenc}
\usepackage{amssymb,amsmath,amsfonts,url}
\usepackage{graphicx,color}
\usepackage{pgf,pgfarrows,pgfnodes,pgfautomata,pgfheaps,pgfshade}
\usepackage[utf8]{inputenc}
\usepackage[french]{babel}
\usepackage{tikz}
\usetikzlibrary{positioning}
\usepackage[all]{xy}
\usepackage{amstext}
\usepackage{xspace}
\usepackage{mdwtab}


\mode<article> {
  \usepackage{fullpage}
  \usepackage{hyperref}
  \setjobnamebeamerversion{advancedcourse.beamer}
}





\newcommand{\alertonce}[1]{\addtocounter{beamerpauses}{-1}\alert<+>{#1}}
\newcommand{\alerttwice}[1]{\addtocounter{beamerpauses}{-1}\alert<+-+(1)>{#1}}

%\newcommand{\subsubsubsection}[1]{\noindent\textbf{#1}}

\newcommand{\cfbox}[1]{\begin{center}\fbox{#1}\end{center}}
\newcommand{\cbox}[1]{\begin{center} #1 \end{center}}

%%%%%%%%%%%%%%%%%%%%%%%%%%%%%%%%%%%%%%%%%%%%%%%%%%%%%%%%%%%%%%%%%%%%%%%%%%%%%%%%%%%%%%%%%

\graphicspath{{Graphics/}}
%\input def_0.tex

%\input{StyleSlides.tex}
\mode<presentation> {
    \usepackage{beamerthemeBoadilla}

  %\beamertemplatetransparentcovereddynamic
  \beamertemplateballitem
}

\usepackage{listings}
\usepackage{xcolor}

\lstset{
	language=Python,
	basicstyle=\ttfamily\small,
	keywordstyle=\color{blue},
	commentstyle=\color{gray},
	stringstyle=\color{green!50!black},
	showstringspaces=false,
	frame=single,
	numbers=none
}
\lstset{
	language=Python,
	basicstyle=\ttfamily\small,
	keywordstyle=\color{blue},
	commentstyle=\color{gray},
	stringstyle=\color{green!50!black},
	showstringspaces=false,
	frame=single,
	numbers=none,
	literate=
	{é}{{\'e}}1
	{è}{{\`e}}1
	{ê}{{\^e}}1
	{ë}{{\"e}}1
	{à}{{\`a}}1
	{â}{{\^a}}1
	{ä}{{\"a}}1
	{ù}{{\`u}}1
	{û}{{\^u}}1
	{ü}{{\"u}}1
	{ô}{{\^o}}1
	{ö}{{\"o}}1
	{î}{{\^i}}1
	{ï}{{\"i}}1
	{ç}{{\c c}}1
}


\title[IN100 Cours 9]{\textbf{Fondements informatiques I}\\
Cours 9: Algorithmes : piles et files}

\author[]{Sorina Ionica \texttt{sorina.ionica@uvsq.fr} \\~\\Bilal FAYE  \texttt{bilal.faye@uvsq.fr} }


\institute{}

\date{}


%\setbeamercolor{normal text}{fg=black}

\begin{document}


\frame{\titlepage}

\begin{frame}[fragile]{File : définition et principe}

Une \textbf{file} (en anglais, \textit{queue}) est une structure
de données dans laquelle :

\begin{itemize}
    \item les éléments sont ajoutés à la fin de la file ;
    \item les éléments sont retirés au début de la file.
\end{itemize}

\begin{exampleblock}{Exemple concret}
Comme dans une file d'attente à la caisse :
la première personne arrivée est la première à être servie.
\end{exampleblock}

\pause

\begin{center}
\textbf{FIFO : First In, First Out}
\end{center}


\begin{tikzpicture}[
    box/.style={draw, minimum width=0.6cm, minimum height=0.6cm, anchor=west},
    node distance=0cm,
    >=stealth
]

% === File initiale ===
\node[box] (a1) at (0,0) {3};
\node[box, right=of a1] (a2) {5};
\node[box, right=of a2] (a3) {7};

% === File après ajout de 4 et 6 ===
\node[box] (b1) at (5,0) {3};
\node[box, right=of b1] (b2) {5};
\node[box, right=of b2] (b3) {7};
\node[box, right=of b3] (b4) {\textcolor{red}{4}};
\node[box, right=of b4] (b5) {\textcolor{red}{6}};

% Flèche d'ajout
\node[draw=none, anchor=west] at (2.5,0.2) {\small{Rajouter 4, 6}};
\draw[->, thick] (a3.east) -- (b1.west);

% === File après un premier retrait ===
\node[box] (c1) at (6.3,-2) {5};
\node[box, right=of c1] (c2) {7};
\node[box, right=of c2] (c3) {4};
\node[box, right=of c3] (c4) {6};

% === File après un deuxième retrait ===
\node[box] (d1) at (10.3,-2) {7};
\node[box, right=of d1] (d2) {4};
\node[box, right=of d2] (d3) {6};

% Flèche de retrait après ajout
\node[draw=none, anchor=west] at (6.5,-0.7) {\small{Enlever}};
\draw[->, thick] (b5.east) .. controls +(0.9,-2.2) and +(-0.8,2.2) .. (c1.west);

% Flèche entre les deux retraits
\node[draw=none, anchor=west] at (8.8,-1.8) {\small{Enlever}};
\draw[->, thick] (c4.east) -- (d1.west);

\end{tikzpicture}


\pause

\textbf{Applications :}
\begin{itemize}
    \item traitement des tâches dans leur ordre d'arrivée ;
    \item gestion des demandes d'impression ;
    \item transmission de données dans un réseau.
\end{itemize}

\end{frame}

\begin{frame}[fragile]{File : les opérations principales}

Pour manipuler une file, on utilise deux opérations :

\begin{itemize}
    \item \textbf{Ajouter :} insérer un élément à la fin de la file.
    \item \textbf{Enlever :} retirer l'élément situé au début de la file.
\end{itemize}

\medskip
\begin{tikzpicture}[
    box/.style={draw, minimum width=0.6cm, minimum height=0.6cm, anchor=west},
    node distance=0cm,
    >=stealth
]

% === File initiale ===
\node[box] (a1) at (0,0) {3};
\node[box, right=of a1] (a2) {5};
\node[box, right=of a2] (a3) {7};

% === File après ajout de 4 et 6 ===
\node[box] (b1) at (5,0) {3};
\node[box, right=of b1] (b2) {5};
\node[box, right=of b2] (b3) {7};
\node[box, right=of b3] (b4) {\textcolor{red}{4}};
\node[box, right=of b4] (b5) {\textcolor{red}{6}};

% Flèche d'ajout
\node[draw=none, anchor=west] at (2.5,0.2) {\small{Rajouter 4, 6}};
\draw[->, thick] (a3.east) -- (b1.west);

% === File après un premier retrait ===
\node[box] (c1) at (6.3,-2) {5};
\node[box, right=of c1] (c2) {7};
\node[box, right=of c2] (c3) {4};
\node[box, right=of c3] (c4) {6};

% === File après un deuxième retrait ===
\node[box] (d1) at (10.3,-2) {7};
\node[box, right=of d1] (d2) {4};
\node[box, right=of d2] (d3) {6};

% Flèche de retrait après ajout
\node[draw=none, anchor=west] at (6.5,-0.7) {\small{Enlever}};
\draw[->, thick] (b5.east) .. controls +(0.9,-2.2) and +(-0.8,2.2) .. (c1.west);

% Flèche entre les deux retraits
\node[draw=none, anchor=west] at (8.8,-1.8) {\small{Enlever}};
\draw[->, thick] (c4.east) -- (d1.west);

\end{tikzpicture}


\pause

\begin{exampleblock}{Exemple}
\begin{lstlisting}[language=Python]
file = [3, 5, 7]
ajouter(file, 4)
ajouter(file, 6)
enlever(file)
enlever(file)
\end{lstlisting}
\end{exampleblock}

\pause

Les éléments sont retirés dans le même ordre que celui
dans lequel ils ont été ajoutés.

\end{frame}
\begin{frame}[fragile]{Ajouter dans une file}

On suppose que la file est représentée par un tableau dynamique.

\medskip

\begin{block}{Principe}
Un nouvel élément est toujours ajouté à la
\textbf{fin de la file}.
\end{block}

\medskip

\begin{exampleblock}{Algorithme : ajouter(f,e)}
\begin{verbatim}
Entrée : une file f et un élément e
Sortie : la file f modifiée

Ajouter e à la fin de f
\end{verbatim}
\end{exampleblock}

\medskip

\begin{center}
\begin{tikzpicture}[
    box/.style={draw, minimum width=0.7cm,
    minimum height=0.65cm},
    >=stealth
]
\node[box] (a) at (0,0) {3};
\node[box,right=0cm of a] (b) {5};
\node[box,right=0cm of b] (c) {7};

\node at (1,-0.65) {\small File initiale};

\node[box] (d) at (4,0) {3};
\node[box,right=0cm of d] (e1) {5};
\node[box,right=0cm of e1] (f1) {7};
\node[box,right=0cm of f1,fill=red!15] (g) {4};

\draw[->,thick] (2,0)--(3.4,0);
\node at (5.1,-0.65) {\small Après ajout de 4};
\end{tikzpicture}
\end{center}

\end{frame}

\begin{frame}[fragile]{Enlever dans une file}

On retire toujours l'élément situé au
\textbf{début de la file}.

\medskip

\textbf{Principe: }
L'élément retiré est le premier élément de la file.
Les autres éléments avancent d'une position.

\medskip

\begin{exampleblock}{Algorithme : enlever(f)}
\begin{verbatim}
Entrée : une file f
Sortie : l'élément retiré et la file modifiée
Si f est vide :
    Signaler que la file est vide
Sinon :
    e <- premier élément de f
    Supprimer le premier élément de f
Retourner e
\end{verbatim}
\end{exampleblock}

\medskip

\begin{center}
\begin{tikzpicture}[
    box/.style={draw, minimum width=0.7cm,
    minimum height=0.65cm},
    >=stealth
]
\node[box,fill=red!15] (a) at (0,0) {3};
\node[box,right=0cm of a] (b) {5};
\node[box,right=0cm of b] (c) {7};

\node at (1,-0.65) {\small Avant retrait};

\draw[->,thick] (3,0)--(4,0);

\node[box] (d) at (4.5,0) {5};
\node[box,right=0cm of d] (e) {7};

\node at (5.2,-0.65) {\small Après retrait};
\node at (0,-0.85) {\small Élément retiré : 3};
\end{tikzpicture}
\end{center}

\end{frame}

\begin{frame}[fragile]{Une file en Python}

En Python, on peut utiliser une liste pour représenter une file.

\begin{columns}[T]

\column{0.49\textwidth}

\begin{block}{Ajouter un élément}
\begin{lstlisting}
def ajouter(f, e):
    f.append(e)
\end{lstlisting}
\end{block}
L'élément est ajouté à la fin de la liste.

\column{0.49\textwidth}

\begin{block}{Enlever un élément}
\begin{lstlisting}
def enlever(f):
    if len(f) > 0:
        return f.pop(0)
    return None
\end{lstlisting}
\end{block}
Le premier élément est retiré et renvoyé.

\end{columns}

\medskip

\begin{exampleblock}{Exemple d'utilisation}
\begin{lstlisting}
f = [3, 5, 7]
ajouter(f, 4)
print(f)          # [3, 5, 7, 4]
e = enlever(f)
print(e)          # 3
print(f)          # [5, 7, 4]
\end{lstlisting}
\end{exampleblock}

\end{frame}
\begin{frame}[fragile]{Suivre les opérations sur une file}

On considère la file initiale :

\begin{center}
\texttt{f = [3, 5, 7]}
\end{center}

\medskip

\begin{exampleblock}{Instructions Python}
\begin{lstlisting}
ajouter(f, 4)
ajouter(f, 6)
e1 = enlever(f)
e2 = enlever(f)
print(f)
print(e1, e2)
\end{lstlisting}
\end{exampleblock}

\pause

\begin{block}{Résultat}
\begin{itemize}
    \item File finale : \texttt{[7, 4, 6]}
    \item Premier élément retiré : \texttt{3}
    \item Deuxième élément retiré : \texttt{5}
\end{itemize}
\end{block}

\medskip

\begin{center}
\textbf{Les premiers éléments ajoutés sont les premiers retirés.}
\end{center}

\end{frame}
\begin{frame}[fragile]{Exemple d'utilisation : gestion des impressions}

Une imprimante reçoit plusieurs demandes d'impression.



\begin{exampleblock}{Situation initiale}
Trois documents sont en attente :
$
f = [\texttt{Doc1},\texttt{Doc2},\texttt{Doc3}]
$
Un nouveau document \texttt{Doc4} arrive.
\end{exampleblock}

\medskip

\begin{block}{Programme Python}
\begin{lstlisting}
f = ["Doc1", "Doc2", "Doc3"]
ajouter(f, "Doc4")
while len(f) > 0:
    document = enlever(f)
    print("Impression :", document)
\end{lstlisting}
\end{block}

\pause

\begin{exampleblock}{Résultat}
\begin{verbatim}
Impression : Doc1
Impression : Doc2
Impression : Doc3
Impression : Doc4
\end{verbatim}
\end{exampleblock}

\end{frame}

% ============================================================
% SLIDE 1 : DEFINITION ET PRINCIPE D'UNE PILE
% ============================================================

\begin{frame}[fragile]{Pile : définition et principe}

Une \textbf{pile} (en anglais, \textit{stack}) est une structure
de données dans laquelle :

\begin{itemize}
    \item les éléments sont ajoutés au sommet de la pile ;
    \item les éléments sont retirés du sommet de la pile.
\end{itemize}

\begin{exampleblock}{Exemple concret}
Comme une pile d'assiettes : on pose une assiette au-dessus
des autres et on retire toujours celle du dessus.
\end{exampleblock}

\begin{center}
\textbf{LIFO : Last In, First Out}
\end{center}

\begin{center}
\begin{tikzpicture}[
    rect/.style={
        draw,
        minimum width=0.55cm,
        minimum height=0.35cm,
        align=center
    },
    >=stealth,
    scale=0.85,
    transform shape
]

% === Pile initiale : 3, 5, 7 ===
\begin{scope}[xshift=0cm]
    \node[rect] (s1_1) at (0,0) {3};
    \node[rect] (s1_2) at (0,0.4) {5};
    \node[rect] (s1_3) at (0,0.8) {7};
\end{scope}

% === Après empiler 4 ===
\begin{scope}[xshift=2.5cm]
    \node[rect] (s2_1) at (0,0) {3};
    \node[rect] (s2_2) at (0,0.4) {5};
    \node[rect] (s2_3) at (0,0.8) {7};
    \node[rect,fill=red!15] (s2_4) at (0,1.2) {4};
\end{scope}

% === Après empiler 6 ===
\begin{scope}[xshift=5cm]
    \node[rect] (s3_1) at (0,0) {3};
    \node[rect] (s3_2) at (0,0.4) {5};
    \node[rect] (s3_3) at (0,0.8) {7};
    \node[rect] (s3_4) at (0,1.2) {4};
    \node[rect,fill=red!15] (s3_5) at (0,1.6) {6};
\end{scope}

% === Après dépiler 6 ===
\begin{scope}[xshift=7.5cm]
    \node[rect] (s4_1) at (0,0) {3};
    \node[rect] (s4_2) at (0,0.4) {5};
    \node[rect] (s4_3) at (0,0.8) {7};
    \node[rect] (s4_4) at (0,1.2) {4};
\end{scope}

% === Après dépiler 4 ===
\begin{scope}[xshift=10cm]
    \node[rect] (s5_1) at (0,0) {3};
    \node[rect] (s5_2) at (0,0.4) {5};
    \node[rect] (s5_3) at (0,0.8) {7};
\end{scope}

% === Flèches entre les piles ===
\draw[->,thick]
    (s1_3.east) .. controls +(right:8mm) and +(left:8mm) ..
    node[midway,above] {\scriptsize Empiler 4}
    (s2_4.west);

\draw[->,thick]
    (s2_4.east) .. controls +(right:8mm) and +(left:8mm) ..
    node[midway,above] {\scriptsize Empiler 6}
    (s3_5.west);

\draw[->,thick]
    (s3_5.east) .. controls +(right:8mm) and +(left:8mm) ..
    node[midway,above] {\scriptsize Dépiler}
    (s4_4.west);

\draw[->,thick]
    (s4_4.east) .. controls +(right:8mm) and +(left:8mm) ..
    node[midway,above] {\scriptsize Dépiler}
    (s5_3.west);

% === Légende ===
\node at (0,-0.45) {\scriptsize Initiale};
\node at (2.5,-0.45) {\scriptsize Après ajout};
\node at (5,-0.45) {\scriptsize Après ajout};
\node at (7.5,-0.45) {\scriptsize Après retrait};
\node at (10,-0.45) {\scriptsize Après retrait};

\end{tikzpicture}
\end{center}

\medskip

\textbf{Applications :}
\begin{itemize}
    \item inversion d'une séquence ;
    \item évaluation d'expressions arithmétiques ;
    \item historique de navigation dans un navigateur.
\end{itemize}

\end{frame}


% ============================================================
% SLIDE 2 : REPRESENTATION ET EMPILER
% ============================================================

\begin{frame}[fragile]{Pile : représentation et empilement}

On représente la pile par un tableau dynamique.

\begin{itemize}
    \item Le sommet correspond à la dernière case du tableau.
    \item Empiler consiste à ajouter un élément à la fin.
\end{itemize}

\medskip

\begin{exampleblock}{Algorithme : empiler(p,e)}
\begin{lstlisting}
Entree : une pile p et un element e
Sortie : la pile p modifiee

Ajouter e a la fin de p
\end{lstlisting}
\end{exampleblock}

\medskip

\begin{center}
\begin{tikzpicture}[
    box/.style={
        draw,
        minimum width=0.7cm,
        minimum height=0.55cm
    },
    >=stealth
]

\node[box] (a) at (0,0) {3};
\node[box] (b) at (0.7,0) {5};
\node[box] (c) at (1.4,0) {7};

\node at (0.7,-0.55) {\small Avant};

\draw[->,thick] (2.2,0)--(3.1,0)
    node[midway,above] {\small Empiler 4};

\node[box] (d) at (3.8,0) {3};
\node[box] (e) at (4.5,0) {5};
\node[box] (f) at (5.2,0) {7};
\node[box,fill=red!15] (g) at (5.9,0) {4};

\node at (4.85,-0.55) {\small Après};

\node at (1.4,0.65) {\small Sommet};
\draw[->] (1.4,0.55)--(1.4,0.3);

\node at (5.9,0.65) {\small Nouveau sommet};
\draw[->] (5.9,0.55)--(5.9,0.3);

\end{tikzpicture}
\end{center}

\end{frame}


% ============================================================
% SLIDE 3 : DEPILER
% ============================================================

\begin{frame}[fragile]{Pile : dépiler}

Dépiler consiste à retirer l'élément situé au sommet de la pile.

\medskip

\begin{exampleblock}{Algorithme : depiler(p)}
\begin{lstlisting}
Entree : une pile p
Sortie : l'element retire

Si p est vide :
    Signaler que la pile est vide
Sinon :
    e <- dernier element de p
    Supprimer le dernier element de p
    Retourner e
\end{lstlisting}
\end{exampleblock}

\medskip

\begin{center}
\begin{tikzpicture}[
    box/.style={
        draw,
        minimum width=0.7cm,
        minimum height=0.55cm
    },
    >=stealth
]

\node[box] (a) at (0,0) {3};
\node[box] (b) at (0.7,0) {5};
\node[box,fill=red!15] (c) at (1.4,0) {7};

\node at (0.7,-0.55) {\small Avant};

\draw[->,thick] (2.2,0)--(3.1,0)
    node[midway,above] {\small Dépiler};

\node[box] (d) at (3.8,0) {3};
\node[box] (e) at (4.5,0) {5};

\node at (4.15,-0.55) {\small Après};

\node at (1.4,0.65) {\small Sommet};
\draw[->] (1.4,0.55)--(1.4,0.3);

\node at (6,0) {\small Élément retiré : 7};

\end{tikzpicture}
\end{center}

\end{frame}


% ============================================================
% SLIDE 4 : IMPLEMENTATION EN PYTHON
% ============================================================

\begin{frame}[fragile]{Une pile en Python}

En Python, on utilise une liste pour représenter une pile.


\begin{columns}[T]

\column{0.48\textwidth}

\begin{block}{Empiler}
\begin{lstlisting}
def empiler(p, e):
    p.append(e)
\end{lstlisting}

Ajoute l'element au sommet.
\end{block}

\column{0.48\textwidth}

\begin{block}{Dépiler}
\begin{lstlisting}
def depiler(p):
    if len(p) > 0:
        return p.pop()
    return None
\end{lstlisting}

Retire et renvoie le sommet.
\end{block}

\end{columns}

\medskip

\begin{exampleblock}{Exemple d'utilisation}
\begin{lstlisting}
p = [3, 5, 7]
empiler(p, 4)
empiler(p, 6)
print(p)          # [3, 5, 7, 4, 6]
depiler(p)
depiler(p)
print(p)          # [3, 5, 7]
\end{lstlisting}
\end{exampleblock}

\end{frame}


% ============================================================
% SLIDE 5 : APPLICATION - INVERSER UNE SEQUENCE
% ============================================================

\begin{frame}[fragile]{Exemple d'utilisation : inverser une séquence}

Une pile permet d'inverser l'ordre des éléments d'un tableau.

\medskip

\begin{block}{Principe}
On retire successivement les éléments du tableau
et on les empile dans une nouvelle pile.
\end{block}

\medskip

\begin{exampleblock}{Programme Python}
\begin{lstlisting}
a = [3, 5, 7, 2]
inverse = []
while len(a) > 0:
    e = depiler(a)
    empiler(inverse, e)
print(inverse)
\end{lstlisting}
\end{exampleblock}

\pause

\begin{center}
\begin{tikzpicture}[
    box/.style={
        draw,
        minimum width=0.7cm,
        minimum height=0.55cm
    },
    >=stealth
]

% Tableau initial
\node[box] (a) at (0,0) {3};
\node[box] (b) at (0.7,0) {5};
\node[box] (c) at (1.4,0) {7};
\node[box,fill=red!15] (d) at (2.1,0) {2};

\node at (1,-0.55) {\small Tableau initial};

\draw[->,thick] (3,0)--(4,0)
    node[midway,above] {\small Inverser};

% Tableau inversé
\node[box,fill=red!15] (e) at (4.8,0) {2};
\node[box] (f) at (5.5,0) {7};
\node[box] (g) at (6.2,0) {5};
\node[box] (h) at (6.9,0) {3};

\node at (5.9,-0.55) {\small Séquence inversée};

\end{tikzpicture}
\end{center}

\end{frame}









\begin{comment}
\begin{frame}[fragile]{File}

Une file (en anglais queue) est une structure de données basée sur le principe que les premiers éléments ajoutés dans la file seront les premiers à en être retirés.


\vspace{0.5 cm}


\begin{tikzpicture}[
    box/.style={draw, minimum width=0.6cm, minimum height=0.6cm, anchor=west},
    node distance=0cm,
    >=stealth
]

% === Initial queue ===
\node[box] (a1) at (0,0) {3};
\node[box, right=of a1] (a2) {5};
\node[box, right=of a2] (a3) {7};

%\node[draw=none, anchor=west] at (0,-0.8) {Initial queue};

% === Queue after enqueue 4 and 6 (same level) ===
\node[box] (b1) at (5,0) {\textcolor{red}{6}};
\node[box, right=of b1] (b2) {\textcolor{red}{4}};
\node[box, right=of b2] (b3) {3};
\node[box, right=of b3] (b4) {5};
\node[box, right=of b4] (b5) {7};


% Arrow from initial queue to enqueue queue
\node[draw=none, anchor=west] at (2.5,0.2) {\small{Rajouter 4, 6}};
\draw[->, thick] (a3.east) -- (b1.west);

% === Queue after two dequeue operations ===
\node[box] (c1) at (6.3,-2) {6};
\node[box, right=of c1] (c2) {4};
\node[box, right=of c2] (c3) {3};
\node[box, right=of c3] (c4) {5};

\node[draw=none, anchor=west] at (8.8,-1.8) {\small{Enlever}};


\node[box] (d1) at (10.3,-2) {6};
\node[box, right=of d1] (d2) {4};
\node[box, right=of d2] (d3) {3};

\draw[->, thick] (c4.east) -- (d1.west);

%\node[draw=none, anchor=west] at (6,-3.8) {After dequeue twice};

% Arrow from enqueue queue to dequeue queue
\node[draw=none, anchor=west] at (6.5,-0.7) {\small{Enlever}};
\draw[->, thick] (b5.east) .. controls +(0.9,-2.2) and +(-0.8,2.2) .. (c1.west);

\end{tikzpicture}

\pause

\vspace{0.3 cm}

Nombreuses applications: 
\begin{itemize}
\item traiter dans l'ordre une liste de transactions, la gestion d'un stock
\item gestion des tâches dans l'ordre d'arrivée par exemple sur un système d'exploitation 
\item transmission des  données dans l'ordre sur un réseau
\end{itemize}


\end{frame}


\begin{frame}[fragile]{File}


\vspace{0.5 cm}

\begin{tikzpicture}[
    box/.style={draw, minimum width=0.6cm, minimum height=0.6cm, anchor=west},
    node distance=0cm,
    >=stealth
]

% === Initial queue ===
\node[box] (a1) at (0,0) {3};
\node[box, right=of a1] (a2) {5};
\node[box, right=of a2] (a3) {7};

%\node[draw=none, anchor=west] at (0,-0.8) {Initial queue};

% === Queue after enqueue 4 and 6 (same level) ===
\node[box] (b1) at (5,0) {\textcolor{red}{6}};
\node[box, right=of b1] (b2) {\textcolor{red}{4}};
\node[box, right=of b2] (b3) {3};
\node[box, right=of b3] (b4) {5};
\node[box, right=of b4] (b5) {7};


% Arrow from initial queue to enqueue queue
\node[draw=none, anchor=west] at (2.5,0.2) {\small{Rajouter 4, 6}};
\draw[->, thick] (a3.east) -- (b1.west);

% === Queue after two dequeue operations ===
\node[box] (c1) at (6,-2) {6};
\node[box, right=of c1] (c2) {4};
\node[box, right=of c2] (c3) {3};
\node[box, right=of c3] (c4) {5};

\node[draw=none, anchor=west] at (8.5,-1.8) {\small{Enlever}};


\node[box] (d1) at (10,-2) {6};
\node[box, right=of d1] (d2) {4};
\node[box, right=of d2] (d3) {3};

\draw[->, thick] (c4.east) -- (d1.west);

%\node[draw=none, anchor=west] at (6,-3.8) {After dequeue twice};

% Arrow from enqueue queue to dequeue queue
\node[draw=none, anchor=west] at (6.5,-0.7) {\small{Enlever}};

\draw[->, thick] (b5.east) .. controls +(0.9,-2.2) and +(-0.8,2.2) .. (c1.west);

\end{tikzpicture}


\vspace{0.3cm}

Pour mettre en place une file nous avons besoin de deux procédures (fonctions en Python): 
\begin{itemize}
\item une pour ajouter en tête de file 
\item une pour enlever de la file. 
\end{itemize}

\vspace{-0.25cm}
\pause

\begin{columns}
\column{5 cm}
\begin{exampleblock}{}
\begin{verbatim}
file=[3,5,7]
ajouter(file,4)
ajouter(file,6)
enlever(file)
enlever(file)
\end{verbatim}
\end{exampleblock}
\end{columns}


\end{frame}

\frame[containsverbatim]{
\frametitle{Ajouter dans le file}

On suppose que notre pile est représentée par un tableau dynamique. 

\medskip
\begin{alertblock}{Procédure ajouter}
\begin{verbatim}
Input: La file f et l'élément e à ajouter
Output : La file modifiée 

l<-longueur(f)
On agrandit f d'une case (c.a.d. on alloue de la mémoire, 
longueur(f)<- longueur(f)+1)
Pour tout i allant de l-1 à 0
     f[i+1]<- f[i]
     
f[0]=e
Return f
\end{verbatim}
\end{alertblock}


}


\frame[containsverbatim]{
\frametitle{Enlever de la file}

On suppose que notre pile est représentée par un tableau dynamique. 

\medskip
\begin{alertblock}{Procédure enlever}
\begin{verbatim}
Input: La file f 
Output : Le dernier élément de la file et  la file modifiée 

l<-longueur(f)

Si (l>0) alors: 
       e<-f[l-1]
       On rétrécit f d'une case (désallouer mémoire, 
       longueur(f)<-l-1)
Return e,f
\end{verbatim}
\end{alertblock}


}


\begin{frame}[fragile]
{Une file en Python}

En Python, on utilisera les listes pour représenter des tableaux dynamiques. 

\begin{exampleblock}{}
\begin{columns}
\column{5 cm}
\begin{verbatim}
def ajouter(f,e):
    l=len(f)
    f.append(0)
    for i in range(l):
   	   f[l-i]=f[l-i-1]
    f[0]=e
\end{verbatim}

\pause

\column{5 cm}
\vspace{-1.5 cm}
\begin{verbatim}
def enlever(f):
    if (len(f)>0):
        return f.pop()
\end{verbatim}
\end{columns}
\end{exampleblock}


\pause

\vspace{0.5 cm}


\vspace{0.2 cm}

La fonction \texttt{pop(i)} enlève l'élément d'indice \texttt{i} dans la liste. \\

\vspace{0.2 cm}

L'appel \texttt{pop()} enlève le dernier élément de la liste. 

\vspace{0.2 cm}
Serait-il possible d'utiliser la fonction \texttt{remove} au lieu de \texttt{pop}?  \\


\end{frame}


\begin{frame}[fragile]
{Une file en Python}



\begin{exampleblock}{}
\begin{columns}
\column{5 cm}
\vspace{-0.275 cm}
\begin{verbatim}
def ajouter(f,e):
    l=len(f)
    f.append(0)
    for i in range(l):
   	   f[l-i]=f[l-i-1]
    f[0]=e
\end{verbatim}

\column{5 cm}
\vspace{-1.5 cm}
\begin{verbatim}
def enlever(f):
    if (len(f)>0):
       return f.pop()
\end{verbatim}
\end{columns}
\end{exampleblock}

\vspace{0.5 cm}

Pour quoi ces fonctions ne renvoient pas \texttt{f} comme demandé? 
\pause

\begin{exampleblock}{}
\begin{verbatim}
f=[3,5,7]
ajouter(f,4)
ajouter(f,6)    
print("voici la file",f)
enlever(f)
print("voici la file",f)
\end{verbatim}
\end{exampleblock}



\end{frame}





\begin{frame}[fragile]{Exemple d'utilisation}

\'Ecrire un programme qui permet de  traiter les  transactions sur un compte bancaire pendant une journée. 
%Les recettes seront ajoutées au montant du compte, les dépenses seront déduites sur le compte. 
A la fin de la journée, le montant restant sur le compte est affiché. 

%Input: Un tableau tab qui contient la liste des transactions de la journée et le montant initial sur le compte 
%Output: Le montant restant sur le compte en fin de journée 

\vspace{0.5 cm}
En Python, on utilisera  la fonction \texttt{enlever}. 


\begin{exampleblock}{}
\begin{verbatim} 
montant=3498
f=[53,-78,198,-23,-243]

while len(f)>0: 
   transaction=enlever(f)
   montant=montant+transaction
   
print("montant restant", montant)    
\end{verbatim}
\end{exampleblock}


\end{frame}

\begin{frame}[fragile]{Pile}

\vspace{-0.5 cm}

Une pile (en anglais stack) est une structure de données  fondée sur le principe "dernier arrivé, premier sorti"  \\
Le  dernier élément ajouté à la pile est le premier à en sortir. 
  
  \vspace{0.8 cm}
  
  
  \begin{tikzpicture}[
  rect/.style={draw, minimum width=0.5cm, minimum height=0.3cm, align=center},
  stack/.style={inner sep=0pt},
  >=stealth
]

% === First stack (Initial) at x=0 ===
\begin{scope}[xshift=0cm]
  % stack from bottom (y=0) to top (y=2*0.8)
  \node[rect] (s1_1) at (0,0) {3};
  \node[rect] (s1_2) at (0,0.4) {5};
  \node[rect] (s1_3) at (0,0.8) {7};
\end{scope}

% === Second stack (after pushes) at x=5cm ===
\begin{scope}[xshift=2.5cm]
  \node[rect] (s2_1) at (0,0) {3};
  \node[rect] (s2_2) at (0,0.4) {5};
  \node[rect] (s2_3) at (0,0.8) {7};
  \node[rect] (s2_4) at (0,1.2) {\textcolor{red}{4}}; % pushed 4
  %\node[rect] (s2_5) at (0,1.6) {6}; % pushed 6 (top)
\end{scope}

% === Third stack (after two pops) at x=10cm ===
\begin{scope}[xshift=5cm]
  \node[rect] (s3_1) at (0,0) {3};
  \node[rect] (s3_2) at (0,0.4) {5};
  \node[rect] (s3_3) at (0,0.8) {7}; % back to original top
   \node[rect] (s3_4) at (0,1.2) {4}; % pushed 4
  \node[rect] (s3_5) at (0,1.6) {\textcolor{red}{6}}; % pushed 6 (top)
\end{scope}

\begin{scope}[xshift=7.5cm]
  \node[rect] (s4_1) at (0,0) {3};
  \node[rect] (s4_2) at (0,0.4) {5};
  \node[rect] (s4_3) at (0,0.8) {7}; % back to original topdépiler
   \node[rect] (s4_4) at (0,1.2) {4}; % pushed 4
\end{scope}

\begin{scope}[xshift=10cm]
  \node[rect] (s5_1) at (0,0) {3};
  \node[rect] (s5_2) at (0,0.4) {5};
  \node[rect] (s5_3) at (0,0.8) {7}; % back to original top

\end{scope}


% === Arrows between stacks ===
% arrow from top of initial stack to top of after-push stack
\draw[->, thick] (s1_3.east) .. controls +(right:8mm) and +(left:8mm) .. (s2_4.west)
  node[midway,above] {\small{empiler 4}};

% arrow from after-push stack to after-pop stack
\draw[->, thick] (s2_4.east) .. controls +(right:8mm) and +(left:8mm) .. (s3_5.west)
  node[midway,above] {\small{empiler 6}};
  
  
\draw[->, thick] (s3_5.east) .. controls +(right:8mm) and +(left:8mm) .. (s4_4.west)
  node[midway,above] {\small{dépiler}};

\draw[->, thick] (s4_4.east) .. controls +(right:8mm) and +(left:8mm) .. (s5_3.west)
  node[midway,above] {\small{dépiler}};    

% optional: small braces to show stack boundaries (visual aid)
%\draw[dashed] (-0.9,-0.15) rectangle (0.9,2.05);
%\draw[dashed] (4.1,-0.15) rectangle (5.9,3.45);
%\draw[dashed] (9.1,-0.15) rectangle (10.9,2.05);

\end{tikzpicture}

  \vspace{0.8 cm}
  
  \pause


Nombreuses applications: 
\begin{itemize}
\item inverser une séquence 
\item évaluer des expressions arithmétiques (à voir plus tard)
\item dans les navigateurs internet pour sauvegarder les pages visitées 
\end{itemize}
 
 
\end{frame}





\frame[containsverbatim]{
\frametitle{Représenter une pile, empiler}

Par un tableau dynamique : sommet de pile = dernière case
\begin{itemize}
\item empiler = ajouter en dernière position
\item dépiler = supprimer l’élément en dernière position
\end{itemize}


\begin{alertblock}{Algorithme pour empiler}
\begin{verbatim} 
Input: Pile p et élément e
Output:  Pile p modifiée 

On agrandit p d'une case (allocation mémoire etc.) 
On rajoute e. 
Renvoyer p
\end{verbatim}
\end{alertblock}

}

\frame[containsverbatim]{
\frametitle{Dépiler}


\begin{alertblock}{Algorithme pour dépiler}
\begin{verbatim} 
Input: Pile p 
Output:  Pile p modifiée et la valeur dépilée

l<-longueur(f)
Si (l>0) alors: 
       e<-f[l-1]
On rétrécit p d'une case (désallouer mémoire etc.)

Renvoyer p,e
\end{verbatim}
\end{alertblock}



}



\frame[containsverbatim]{
\frametitle{Une pile en Python}

On utilisera une liste pour représenter la pile. 


\begin{exampleblock}{}
\begin{columns}
\column{5 cm}
\vspace{-0.55 cm}
\begin{verbatim} 
def empiler(p,e):
     p.append(e)
\end{verbatim}

\column{5 cm}

\vspace{-0.2 cm}
\begin{verbatim} 
def depiler(p):
   if len(p)>0:
       return p.pop()
\end{verbatim}
\end{columns}
\end{exampleblock}

\begin{exampleblock}{}

\begin{verbatim} 
p=[3,5,7]
empiler(p,4)
empiler(p,6)
print("voici l'état de la pile", p) #p=[3,5,7,4,6]
depiler(p)
depiler(p)
print("voici l'état de la pile", p) #p=[3,5,7]
\end{verbatim}

\end{exampleblock}

}





\begin{frame}[fragile]{Exemple de programme utilisant une pile}


Inverser un tableau ou une chaîne de caractères en utilisant une pile. \\

En Python, on utilise les fonctions \texttt{empiler} et \texttt{depiler}. 

\vspace{0.5 cm}

\begin{exampleblock}{En Python}
\begin{verbatim} 
a = [3, 5, 7, 2]
inverse = []

while (len(a)>0):
        empiler(inverse, depiler(a))

print(inverse) 
\end{verbatim}
\end{exampleblock}

\vspace{0.5 cm}


\end{frame}


\end{comment}







\end{document}
